%! AP Statistics — Unit 4, Lesson 4.7 — Guided Notes (Answer Key)
\documentclass[10pt]{article}
\usepackage{apstats-article}
\usepackage{apstats-key}
\newcommand{\TallMath}[1]{$\displaystyle #1\rule[-1.4em]{0pt}{3.2em}$}

\begin{document}

\pageheader{Unit 4, Lesson 4.7}{Guided Notes --- Answer Key}
\namedateperiod

\vspace{0.12in}

\begin{objectivebox}
\small By the end of this lesson, I will be able to\dots\\[4pt]
\ans{Represent a probability distribution for a discrete random variable as a table or histogram and verify validity.}\\[1pt]
\ans{Build a cumulative probability distribution and use it to answer probability questions.}\\[1pt]
\ans{Interpret a probability distribution by describing shape, center, and spread in context.}
\end{objectivebox}

\vspace{0.1in}

\begin{vocabbox}
\termblanklong{Random variable} \ans{A variable whose value is a numerical outcome of a random process.}
\termblanklong{Discrete random variable} \ans{A random variable that takes a countable number of values, each with an associated probability.}
\termblanklong{Probability distribution} \ans{A table, graph, or function giving $P(X=x)$ for every possible value of $X$.}
\termblanklong{Probability histogram} \ans{A bar graph with values of $X$ on the horizontal axis and $P(X=x)$ on the vertical axis.}
\termblanklong{Cumulative probability distribution} \ans{A table or function giving $P(X \le x)$ for each value of $X$.}
\end{vocabbox}

\vspace{0.1in}

\begin{hookbox}
\small\textbf{Number of pets owned per household:}
\begin{center}
\renewcommand{\arraystretch}{1.3}
\begin{tabular}{@{}ccccc@{}}
Pets ($x$) & 0 & 1 & 2 & 3 \\
\hline
Probability & 0.30 & 0.40 & 0.20 & 0.10 \\
\end{tabular}
\end{center}
\textbf{What is $P(X \le 1)$?} \quad \ans{$0.30 + 0.40 = 0.70$} \quad
\textbf{What must the probabilities sum to?} \quad \ans{1}
\end{hookbox}

\vspace{0.1in}

\begin{notesbox}{Random Variables and Probability Distributions}
\small

\textbf{Definition:} A \textbf{random variable} is a variable whose value is a
\ans{numerical outcome} of a random process.

\vspace{8pt}
A \textbf{discrete random variable} can take a \ans{countable} number of values.
Each value has a probability associated with it.

\vspace{10pt}
\textbf{Valid Probability Distribution --- Two Conditions:}
\begin{enumerate}[leftmargin=2em, itemsep=4pt]
  \item Each probability: \ans{$P(X=x) \ge 0$ (between 0 and 1)}
  \item Sum of all probabilities: \ans{$\sum P(X=x) = 1$}
\end{enumerate}

\vspace{10pt}
\textbf{Example:} $X$ = number of siblings for a randomly chosen student.

\vspace{4pt}
\renewcommand{\arraystretch}{1.7}
\begin{tabularx}{\linewidth}{@{} l X X X X X @{}}
$x$ & 0 & 1 & 2 & 3 & 4 \\
\hline
$P(X=x)$ & \ans{0.15} & \ans{0.40} & \ans{0.30} & \ans{0.10} & \ans{0.05} \\
\end{tabularx}

\vspace{8pt}
Check conditions: \quad (1) \ans{All values $\ge 0$} \quad (2) \ans{$0.15+0.40+0.30+0.10+0.05 = 1.00$ \checkmark}

\end{notesbox}

\vspace{0.1in}

\begin{notesbox}{Cumulative Probability Distribution}
\small

$P(X \le x)$ = the probability that $X$ is \ans{less than or equal to $x$}.

\vspace{8pt}
\textbf{Example:} Use the siblings distribution above.

\renewcommand{\arraystretch}{1.7}
\begin{tabularx}{\linewidth}{@{} l X X X X X @{}}
$x$ & 0 & 1 & 2 & 3 & 4 \\
\hline
$P(X=x)$ & 0.15 & 0.40 & 0.30 & 0.10 & 0.05 \\
$P(X \le x)$ & \ans{0.15} & \ans{0.55} & \ans{0.85} & \ans{0.95} & \ans{1.00} \\
\end{tabularx}

\vspace{8pt}
$P(X \ge 3) = $ \ans{$0.10 + 0.05 = 0.15$}

\vspace{8pt}
Interpret $P(X \le 2) = 0.85$ in context:\\
\ansline{There is an 85\% probability that a randomly chosen student has 2 or fewer siblings.}
\end{notesbox}

\vspace{0.1in}

\begin{notesbox}{Interpreting a Probability Distribution}
\small

When you describe a probability distribution, discuss \textbf{three things}:

\vspace{6pt}
\renewcommand{\arraystretch}{1.8}
\begin{tabularx}{\linewidth}{@{} p{2.2cm} X @{}}
\textbf{Shape} & \ans{Skewed left, skewed right, roughly symmetric, or uniform} \\
\textbf{Center} & Where are the probabilities \ans{concentrated (most likely values)}? \\
\textbf{Spread} & How \ans{spread out or concentrated} are the values? \\
\end{tabularx}

\vspace{8pt}
\textbf{Interpretation of the siblings distribution in context:}\\
\ansline{The distribution is right-skewed. Most students have 1 or 2 siblings (70\% combined),}\\
\ansline{and few students have 3 or 4 siblings. The most likely outcome is 1 sibling.}
\end{notesbox}

\vspace{0.1in}

\begin{practicebox}
\small
$X$ = number of cars owned per household.
\begin{center}
\renewcommand{\arraystretch}{1.3}
\begin{tabular}{@{}ccccc@{}}
$x$ & 0 & 1 & 2 & 3 \\
\hline
$P(X=x)$ & 0.10 & 0.40 & 0.35 & 0.15 \\
\end{tabular}
\end{center}
\begin{enumerate}[leftmargin=1.5em, itemsep=8pt]
  \item \ans{$0.10+0.40+0.35+0.15=1.00$ and all values $\ge 0$. Valid. \checkmark}
  \item $P(X = 2) = $ \ans{$0.35$}
  \item $P(X \le 1) = $ \ans{$0.10+0.40=0.50$}
  \item $P(X \ge 2) = $ \ans{$0.35+0.15=0.50$}
  \item \ansline{There is a 50\% probability that a randomly chosen household owns 2 or more cars.}
\end{enumerate}
\end{practicebox}

\end{document}
